\documentclass[11pt,a4paper]{article}
\usepackage[utf8]{inputenc}
\usepackage[T1]{fontenc}
\usepackage[a4paper,margin=2.2cm]{geometry}
\usepackage{enumitem}
\usepackage{hyperref}
\hypersetup{colorlinks=true,linkcolor=blue,urlcolor=blue}
\title{MedCareSO\\Guide associé : sécurité, licences et livraison client}
\author{Procédure interne MedCareSO}
\date{\today}
\begin{document}
\maketitle

\section*{Règle la plus importante}
La clé privée \texttt{medcareso-license-private.pem} permet de fabriquer une
licence valide. Elle est l'équivalent du tampon de l'entreprise. Elle ne doit
jamais être envoyée à un client, ajoutée à GitHub, copiée dans le logiciel ou
placée dans le dossier \texttt{cabinet}. Une clé publique peut être distribuée;
une clé privée doit rester dans un coffre chiffré accessible seulement aux
associés autorisés.

\section{Ce que fait la protection}
\begin{itemize}
\item Le client reçoit un fichier de licence JSON signé, et non un code texte.
\item La licence est liée à l'identifiant unique du PC du cabinet.
\item L'application vérifie la signature Ed25519, la date d'expiration et
l'identifiant du PC au démarrage puis régulièrement.
\item Une version distribuée vérifie également l'intégrité de fichiers
critiques et masque les outils de développement usuels.
\item Ces mesures découragent la copie et la modification simple, mais aucun
logiciel hors ligne ne peut être inviolable à 100\%.
\end{itemize}

\section{Première configuration, une seule fois}
Les clés ont déjà été créées par MedCareSO. Si la paire doit être créée sur un
nouveau poste développeur, utiliser :
\begin{verbatim}
npm run license:generate-keypair -- --out-dir /dossier-prive-chiffre/medcareso-keys
\end{verbatim}
La clé publique doit être intégrée avant un build client :
\begin{verbatim}
npm run license:install-public-key -- --key /dossier-prive-chiffre/medcareso-keys/medcareso-license-public.pem
\end{verbatim}
Ne recréez pas une paire de clés sans décision commune : toutes les anciennes
licences deviendraient invalides si la clé publique est remplacée.

\section{Créer une version à livrer}
Depuis le dossier du projet, définir le chemin de la clé privée puis construire
la version protégée :
\begin{verbatim}
export MEDCARESO_LICENSE_PRIVATE_KEY_FILE=/dossier-prive-chiffre/medcareso-keys/medcareso-license-private.pem
npm run build:secure
\end{verbatim}
Le résultat est dans \texttt{dist/}. Ne distribuez pas un build normal si une
version sécurisée est demandée. Testez toujours la version livrée sur un poste
ou un profil utilisateur propre.

\section{Activer un nouveau cabinet}
\begin{enumerate}[label=\arabic*.]
\item Installez MedCareSO chez le cabinet.
\item Ouvrez l'écran d'activation et récupérez l'\emph{Identifiant de ce poste}.
\item Gardez cet identifiant avec le nom du cabinet dans votre dossier interne.
\item Créez le fichier de licence, par exemple pour une année :
\begin{verbatim}
npm run license:sign -- \
 --private-key /dossier-prive-chiffre/medcareso-keys/medcareso-license-private.pem \
 --machine-id IDENTIFIANT_DU_POSTE \
 --cabinet "Cabinet Dr. Nom" \
 --expires 2027-07-18 \
 --features standard \
 --out /dossier-sortie/Cabinet-Dr-Nom.medcareso.json
\end{verbatim}
\item Envoyez seulement ce fichier \texttt{.medcareso.json} au cabinet.
\item Le cabinet choisit le fichier dans MedCareSO puis clique sur
\emph{Activer la licence}.
\end{enumerate}

\section{Renouvellement et changement de PC}
Pour renouveler, créez un nouveau fichier signé avec la nouvelle date. Pour un
nouveau PC, demandez le nouvel identifiant de poste et créez une nouvelle
licence : l'ancienne licence n'est pas valable sur une autre machine.

\section{Travail à deux et GitHub}
\begin{itemize}
\item Partagez le dépôt source et les guides via GitHub.
\item Partagez la clé privée seulement par un coffre chiffré commun (Bitwarden,
1Password, coffre VeraCrypt ou support chiffré), jamais par GitHub, WhatsApp ou
email en clair.
\item Vérifiez \texttt{git status} avant chaque push. Les dossiers
\texttt{developer-keys/} et les fichiers \texttt{*.medcareso.json} sont ignorés.
\item Ne publiez jamais un jeton GitHub dans le terminal, dans une capture ou
dans une conversation. Révoquez immédiatement tout jeton exposé.
\end{itemize}

\section*{Checklist avant livraison}
\begin{enumerate}[label={[\ ]}]
\item Build créé avec \texttt{npm run build:secure}.
\item Licence signée pour le bon identifiant de poste.
\item Fichier de licence testé sur le PC du cabinet.
\item Aucun fichier privé ou jeton ajouté au dépôt.
\item Identifiants admin par défaut changés chez le client.
\item Sauvegardes PostgreSQL protégées et accès réseau limité.
\end{enumerate}
\end{document}
